\section{Ex10: Association Rules -- Chiến lược sinh ứng viên}

\textbf{Đề bài:} Cho tập các 3-itemsets thường xuyên: $\{1,2,3\}, \{1,2,4\}, \{1,2,5\}, \{1,3,4\}, \{1,3,5\}, \{2,3,4\}, \{2,3,5\}, \{3,4,5\}$. Giả sử tập dữ liệu có 5 mặt hàng: $\{1,2,3,4,5\}$.

\subsection*{(a) Liệt kê tất cả các 4-itemsets ứng viên sinh ra bằng chiến lược gộp $F_{k-1} \times F_1$}
Chiến lược: Mỗi 3-itemset sẽ được gộp với từng item đơn lẻ trong $F_1$ (từ 1 đến 5) nếu tạo ra tập 4 phần tử mới.

Sinh ứng viên từ từng 3-itemset:
\begin{itemize}
    \item Từ $\{1,2,3\}$: thêm $4 \rightarrow \{1,2,3,4\}$, thêm $5 \rightarrow \{1,2,3,5\}$
    \item Từ $\{1,2,4\}$: thêm $3 \rightarrow \{1,2,3,4\}$, thêm $5 \rightarrow \{1,2,4,5\}$
    \item Từ $\{1,2,5\}$: thêm $3 \rightarrow \{1,2,3,5\}$, thêm $4 \rightarrow \{1,2,4,5\}$
    \item Từ $\{1,3,4\}$: thêm $2 \rightarrow \{1,2,3,4\}$, thêm $5 \rightarrow \{1,3,4,5\}$
    \item Từ $\{1,3,5\}$: thêm $2 \rightarrow \{1,2,3,5\}$, thêm $4 \rightarrow \{1,3,4,5\}$
    \item Từ $\{2,3,4\}$: thêm $1 \rightarrow \{1,2,3,4\}$, thêm $5 \rightarrow \{2,3,4,5\}$
    \item Từ $\{2,3,5\}$: thêm $1 \rightarrow \{1,2,3,5\}$, thêm $4 \rightarrow \{2,3,4,5\}$
    \item Từ $\{3,4,5\}$: thêm $1 \rightarrow \{1,3,4,5\}$, thêm $2 \rightarrow \{2,3,4,5\}$
\end{itemize}
\textbf{Danh sách các ứng viên 4-itemset duy nhất:} $\{1,2,3,4\}, \{1,2,3,5\}, \{1,2,4,5\}, \{1,3,4,5\}, \{2,3,4,5\}$.

\subsection*{(b) Chiến lược gộp Apriori ($F_{k-1} \times F_{k-1}$)}
Theo thuật toán Apriori, hai tập $(k-1)$-itemsets chỉ được gộp với nhau nếu chúng có cùng $k-2$ phần tử đầu tiên. \\
Danh sách các 3-itemsets đã sắp xếp: $123, 124, 125, 134, 135, 234, 235, 345$.

Thực hiện Apriori join:
\begin{itemize}
    \item $123 \text{ join } 124 \rightarrow \{1, 2, 3, 4\}$
    \item $123 \text{ join } 125 \rightarrow \{1, 2, 3, 5\}$
    \item $124 \text{ join } 125 \rightarrow \{1, 2, 4, 5\}$
    \item $134 \text{ join } 135 \rightarrow \{1, 3, 4, 5\}$
    \item $234 \text{ join } 235 \rightarrow \{2, 3, 4, 5\}$
\end{itemize}
\textbf{Kết quả ứng viên Apriori:} $\{1, 2, 3, 4\}, \{1, 2, 3, 5\}, \{1, 2, 4, 5\}, \{1, 3, 4, 5\}, \{2, 3, 4, 5\}$. \\
$\Rightarrow$ \textbf{Đáp án (b):} Trùng với kết quả phần (a).

\subsection*{(c) 4-itemsets sống sót sau bước Apriori Pruning}
Nguyên tắc: Một 4-itemset chỉ được giữ lại nếu mọi tập con 3 phần tử (3-subset) của nó đều nằm trong danh sách các tập mục thường xuyên đã cho.
\begin{enumerate}
    \item $\{1, 2, 3, 4\}$: Các tập con: $123, 124, 134, 234$ (đều phổ biến) $\rightarrow$ \textbf{Giữ lại}.
    \item $\{1, 2, 3, 5\}$: Các tập con: $123, 125, 135, 235$ (đều phổ biến) $\rightarrow$ \textbf{Giữ lại}.
    \item $\{1, 2, 4, 5\}$: Các tập con: $124, 125, 145 \ (\times), 245 \ (\times) \rightarrow$ \textbf{Loại}.
    \item $\{1, 3, 4, 5\}$: Các tập con: $134, 135, 145 \ (\times), 345 \rightarrow$ \textbf{Loại}.
    \item $\{2, 3, 4, 5\}$: Các tập con: $234, 235, 245 \ (\times), 345 \rightarrow$ \textbf{Loại}.
\end{enumerate}
\textbf{Các 4-itemset sống sót cuối cùng:} $\{1, 2, 3, 4\}, \{1, 2, 3, 5\}$. \\
$\Rightarrow$ \textbf{Đáp án (c):} $\{1, 2, 3, 4\}, \{1, 2, 3, 5\}$.

\begin{table}[H]
    \centering
    \begin{tabular}{|c|l|l|}
        \hline
        \textbf{Giai đoạn} & \textbf{Chiến lược} & \textbf{Các tập mục (Itemsets)} \\
        \hline
        (a) & Ứng viên $F_{k-1} \times F_1$ & $\{1,2,3,4\}, \{1,2,3,5\}, \{1,2,4,5\}, \{1,3,4,5\}, \{2,3,4,5\}$ \\ \hline
        (b) & Ứng viên Apriori (Join) & $\{1,2,3,4\}, \{1,2,3,5\}, \{1,2,4,5\}, \{1,3,4,5\}, \{2,3,4,5\}$ \\ \hline
        (c) & Sau khi Prune (Cắt tỉa) & $\{1,2,3,4\}, \{1,2,3,5\}$ \\ \hline
    \end{tabular}
    \caption{Tổng hợp kết quả Ex10}
\end{table}

\textbf{Nhận xét:}
\begin{itemize}
    \item Bước gộp Apriori Join hiệu quả hơn chiến lược $F_{k-1} \times F_1$ thông thường vì nó tận dụng tri thức về thứ tự các mục để giảm thiểu số lượng ứng viên sinh ra ngay từ đầu.
    \item Bước Pruning là cực kỳ quan trọng; nó loại bỏ các tập mục như $\{1,2,4,5\}, \{1,3,4,5\}$ và $\{2,3,4,5\}$ vì chúng chứa các tập con không phổ biến (ví dụ: $\{1,4,5\}$ hoặc $\{2,4,5\}$), giúp tối ưu hóa hiệu suất tính toán đáng kể.
\end{itemize}